\documentclass[9pt]{beamer}

%-------------------------------------------------
%   THEMES & PACKAGES
%-------------------------------------------------
\usetheme[progressbar=frametitle]{metropolis}
\usepackage[belowskip=-15pt]{caption}
\usepackage{graphicx}

%-------------------------------------------------
%   Settings
%-------------------------------------------------
\setbeamertemplate{footline}[text line]{%
    \parbox{\linewidth}{\vspace*{-8pt}\insertshorttitle\hfill\insertshortsubtitle\hfill\insertshortauthor\hfill\insertpagenumber}}
\setbeamertemplate{navigation symbols}{}

%-------------------------------------------------
%   TITLE
%-------------------------------------------------
\title{Neural Networks}
\subtitle{ESN Paper Review}
\date{Lecture date: Jan 16, 2017}
\author{Minh Nguyen}
\titlegraphic{\hfill\includegraphics[height=0.7cm]{../../images/h-brs-logo.jpg}}

%-------------------------------------------------
%   COMMANDS
%-------------------------------------------------
\newcommand{\picEqHereWidth}[2] { %
    \begin{figure}[htp] 
        \centering
        \includegraphics[width=#2]{#1}
    \end{figure}
}
\newcommand{\picHereWidth}[4] { %
    \begin{figure}[htp] %
        \centering
        \includegraphics[width=#4]{#1} %
        \caption{#2} %
        \label{#3}
    \end{figure} %
}
%-------------------------------------------------
%   BEGIN
%-------------------------------------------------
\begin{document}

%-------------------------------------------------
\maketitle

%-------------------------------------------------
\begin{frame}{Introduction}
    \begin{alertblock}{Title}
        Opening the Black Box: Low-Dimensional Dynamics in High-Dimensional Recurrent Neural Networks
    \end{alertblock}
    \begin{alertblock}{Author}
        \emph{David Sussillo and Omri Barak}
    \end{alertblock}
    \begin{alertblock}{Abstract}
        \begin{itemize}
            \item Explore hypothesis that linearized dynamics around stable and unstable fixed points reveal how RNNs implement their computations
            \item Explore utility of linearization around points not which are not true fixed points but of slow movement
            \item Present simple optimization technique to find fixed and slow points of RNNs' dynamics
            \item Showcase the technique on three high-dimensional RNN examples: 3-bit flip-flop device, input-dependent sine wave generator, \& two-point moving average.
        \end{itemize}
    \end{alertblock}
\end{frame}

%-------------------------------------------------
\begin{frame}{Fixed Points}
    \begin{alertblock}{Definitions and properties}
        \begin{itemize}
            \item Fixed points are common anchors to start an analysis of a nonlinear dynamical system (NLDS)
            \item RNNs can be viewed as NLDSs to learn how they implement their target functions
            \item Fixed points can be stable or unstable (starting in vicinity converges towards or diverges away)
            \item Saddle points with with both stable and unstable modes allow better understanding of RNNs and can be found using the introduced optimization technique
        \end{itemize}
    \end{alertblock}
\end{frame}

%-------------------------------------------------
\begin{frame}{Optimization technique}
    \begin{alertblock}{Linear approximations}
        Consider system of first order differential equations $\mathbf{\dot{x} = F(x)}$. Linearization of the Taylor expansion can be done for $\mathbf{x^*}$ and perturbations $\mathbf{\delta x \equiv x - x^* }$ s.t.:
        \picEqHereWidth{../images/sussilo13-eq3_4.png}{0.3\linewidth}
        Define auxiliary scaler function:
        \[ q(\mathbf{x}) = \frac{1}{2} |\mathbf{F(x)}|^2 \]
        Minimizing $q(\mathbf{x})$ can give candidate regions for linearization. Fixed points are $\mathbf{x^*}$ s.t $q(\mathbf{x}) = 0$
    \end{alertblock}
\end{frame}

%-------------------------------------------------
\begin{frame}{Examples -- 2D}
    \picHereWidth{../images/sussilo13-fig1_2d.png}{Behavior of minimizing $q$ relative to a simple 2D system}{fig:2d}{0.4\linewidth}
    \begin{itemize}
        \item The system (trajectories in black lines): $\dot{x_1} = (1 - x_1^2)x_2,\dot{x_2} = x_1/2 - x_2$
        \item Black `x': 2 attractors; green `x': saddle point
        \item Thick red lines: gradient descent on $q$
        \item Trajectories only lead to fixed points, whereas minimization of $q$ can converge at saddle point
    \end{itemize}
\end{frame}

%-------------------------------------------------
\begin{frame}{Examples -- RNNs}
    \begin{alertblock}{Formulation}
        \picEqHereWidth{../images/sussilo13-eq10_11.png}{0.37\linewidth}
        \begin{itemize}
            \item $\mathbf{r} = h(\mathbf{x})$ are firing rates (element-wise application of nonlinear $h$ to $\mathbf{x}$)
            \item Recurrence defined via $\mathbf{J}$, $I$-dimensional inputs $\mathbf{u}$ are connected via $\mathbf{B}$
            \item Jacobian of network:
            \picEqHereWidth{../images/sussilo13-eq12.png}{0.25\linewidth}
            where $\delta_ij = 1$ if $i = j$ and $0$ otherwise.
        \end{itemize}
    \end{alertblock}
\end{frame}

%-------------------------------------------------
\begin{frame}{Examples -- 3-Bit Flip-Flop}
    \picHereWidth{../images/sussilo13-fig2_flipflop.png}{Left: example input-output pairs of the 3-bit flip-flop. Right: ESN used to learn the flip-flop task}{fig:flipflop_model}{0.55\linewidth}
    \begin{itemize}
        \item Output should transition to or stay at $+1$ ($-1$) when receive input pulse $+1$ ($-1$)
        \item All possible transitions between memory states are spanned using the network inputs, 600 of these trajectories are randomly selected for $q$ optimization.
        \item In figure \ref{fig:flipflop_result} (left) network state $\mathbf{x}(t)$ is projected on first 3 principle components of the network activations. Right panel visualizes transition $(-1,-1,-1)$ to $( 1, -1, -1)$ projected on (PC3, relevant input).
    \end{itemize}
\end{frame}

%-------------------------------------------------
\begin{frame}{Examples -- 3-Bit Flip-Flop}
    \picHereWidth{../images/sussilo13-fig3_flipflop.png}{Left: low-dim phase representation of flip-flop task. Right: demonstration that saddle point mediates transitions between attractors.}{fig:flipflop_result}{0.78\linewidth}
    \begin{itemize}
        \item Black `x': 8 memory states; green `x': saddle fixed points with one unstable dimension; pink `x': fixed points with 4 unstable dimensions.
        \item Blue lines: 24 1-bit transitions; thick red lines: denotes dimensions of instability; thin red lines: trajectories starting from unstable dimension.
        \item For right panel transition input pulse amplitude is varied in small increments from 0 to 1. Blue line is the normal transition and cyan lines correspond with intermediate amplitudes.
    \end{itemize}
\end{frame}

%-------------------------------------------------
\begin{frame}{Examples -- Input-dependent sine wave generator}
    \picHereWidth{../images/sussilo13-fig4_sinegen.png}{Sine wave generator task. \textbf{A}: network is to generate unity amplitude sine wave with frequency given by a static input. \textbf{B}: trajectories (blue), fixed points (green) and eigenvectors of unstable plane around fixed points, projected onto first 3 PCs. \textbf{C}: target frequency compared with that predicted by fixed points.}{fig:sine_gen}{0.7\linewidth}
\end{frame}

%-------------------------------------------------
\begin{frame}{Points of slow movement -- 2-point moving average}
    \picHereWidth{../images/sussilo13-fig7_movingavg.png}{2D approximate plane attractor in 2-point moving average task. \textbf{A}: slow points form a 2D manifold. \textbf{B}: dynamical picture of 2D manifold of slow points.}{fig:sine_gen}{0.68\linewidth}
    \begin{itemize}
        \item Task: predict average of last 2 analog input pulses, which are generated at random times.
        \item gray `x': 2 fixed points; orange: trajectory for computing average with a new input.
        \item green lines: stable direction from fixed points; red lines: unstable directions; blue lines: trajectories from slow points (where network starts).
    \end{itemize}

\end{frame}

%-------------------------------------------------
\begin{frame}{Points of slow movement -- 2-point moving average}
    \picHereWidth{../images/sussilo13-fig8_movingavg.png}{Transition dynamics of RNN vs LTVDS. \textbf{A}: separated trials of RNNs (red) and LTVDS (dotted). \textbf{B}: same trials in space of first 3 PCs; black dots: starting slow points. \textbf{C, D}: transition of values from input pulse to previous 2 inputs; blue are RNN transitions, red are LTVDS.}{fig:sine_gen}{0.5\linewidth}
    \begin{itemize}
        \item Replace the full nonlinear system with a succession of linear systems defined by
        the slow points
        \item Linear Time-Varying Dynamical System (LTVDS): dynamics at each time point are linear dynamics induced by closest slow points.
    \end{itemize}

\end{frame}

%-------------------------------------------------
%   END
%-------------------------------------------------
\end{document}
